% Pursuing Stacks (cotes 134-2 to 134-8): archive cote 134-7, site page 92 (PDF page 93), typescript page 534.
% typescript. Draft compared: ps6.tex:4629-4690.
% AI-transcribed by Claude Opus 5.5 (claude-opus-5-5, adaptive thinking, effort high), Message Batches API, batch ps4 (msgbatch_012Aoc1NuskATPsyBeKSYK7J), 2026-11-12; not yet checked by eye.
\scanpage{134-7/92}
\mspage{534}

description gives rise at once to the familiar ``tower'' structure for
variable n
\[ (\over{\unsure{4}}{5}) \qquad X_{n+1} \to X_n \to \dots \to X_1 \to X_o \quad ( \unsure{=} e ) \quad . \]
Intuitively, the maps in (\over{\unsure{4}}{5}) are viewed as being (surjective) fibrations
between (connected) ``spaces'', X being viewed as a kind of inverse limit
of the $X_n$'s . Dually, we would expect to get a sequence of inclusions
\[ \hand{(*)} \qquad {}_oX \hookrightarrow {}_1X \hookrightarrow \dots \hookrightarrow {}_nX \hookrightarrow {}_{(n+1)}X \hookrightarrow \cdots , \]
with X appearing as \xout{the} a kind of direct limit. In the set up of ordina-%
ry homotopy types, modelized by semisimplicial sets, one \deleted{is}\inserted{will} think\deleted{ing}
at once of the filtration by skeleta - which however is'nt quite the
right thing surely, because if we dualize the familiar Cartan-Serre
requirement on (2) (namely that it induces an isomorphism for $\hand{\pi}_i$ for
$i \leqslant n$), we see that the ``inclusion''
\[ {}_nX \hookrightarrow X \]
should induce an isomorphism\xout{s} for $H_i$ for $i \leqslant n$ , which is'nt quite
true for the skeletal filtration (it is OK for $i \leqslant n-1$ only); the condi-%
tion
\[ H_i(X_n) = 0 \quad \text{for} \quad i > n \]
is OK though for this filtration. So the next idea would be to modify
a little the n-skeleton ${}_nX$ , to straighten things out .\xout{xxxthe analogy}
I played around some along these lines, and after some initial optimism
came to the feeling that there does not exist (in the discrete nor in the
schematic set-up) such an increasing \inserted{canonical} filtration of a homotopy type. I
did'nt make any formal statement and proof for this (in the set-up of
ordinary homotopy types say), however in the process of playing around
it became soon clear that in various ways, \uline{there is} \deleted{an}\inserted{\uline{some}} \uline{essential dissy-}%
\uline{metry} between the seemingly ``dual'' situations, when trying to get the
two types of ``filtrations'' of the object X. O\over{\unsure{n}}{n}e dissymetry occurs alrea-%
dy in the very definition of the subcategory (3), and of the correspon-%
ding ``dual'' \xout{definition} subcategory
\[ (\over{5}{6}) \qquad {}_{(\leqslant n)}Hot_o(k) \]
of objects satisfying
\[ H_i(Y) = 0 \quad \text{for} \quad i > n \ . \]
Namely, both subcategories (3) and (\over{5}{6}) are defined in terms of the \uline{same}
subcategory of $Hotab_o(k)$, namely (\over{\unsure{3}}{4}), as the inverse image of the

\margin{92}
